% Revised Section 4.3 — Acuity-Weighted Class-Conditional Risk Control
%
% Replaces subsections 4.3.1–4.3.4 in manuscript_draft.
%
% Changes from the submitted draft:
%   4.3.1 — removed "quantile" from "stable per-class quantile estimate";
%            added K=35 joint-model statement.
%   4.3.2 — added paragraph explaining AFIB and AFLT are Critical because
%            PTB-XL has no VT/VF; explicit "most-severe-available" caveat.
%   4.3.3 — replaced split-conformal / quantile-q_k language with LTT
%            binomial description; corrected prediction-set definition from
%            {k : 1-f(x)_k <= q_k} to {k : f(x)_k >= lambda_c};
%            replaced "quantile step" sample-floor with LTT certification
%            floor as primary criterion; fixed "Three classes" -> "Two classes"
%            after AFLT removal (AFLT n=7 < 10, never exceeded the naive floor).
%   4.3.4 — unchanged; reproduced for completeness.

% ─────────────────────────────────────────────────────────────────────────────
\subsubsection{Class granularity}
\label{sec:granularity}

We train a single 35-class model that jointly covers the subdiagnostic (23-class)
and rhythm (12-class) label sets, operating at a finer granularity than the coarser
5-class superdiagnostic task or the full 71-statement level.
The superdiagnostic level conflates conditions of substantially different clinical
urgency within a single class (e.g., the CD superclass contains both complete AV
block and isolated incomplete bundle branch block); the full statement level, while
more granular, produces calibration sets too sparse for several rare but high-acuity
statements to support a stable per-class calibration.
The subdiagnostic and rhythm granularities provide classes that are both clinically
distinguishable by acuity and adequately represented in the PTB-XL calibration folds.

% ─────────────────────────────────────────────────────────────────────────────
\subsubsection{Acuity tier definition}
\label{sec:tiers}

Each class is assigned to one of three ordinal acuity tiers---Critical, Important,
Benign---together with a fourth, Excluded, category for baseline/negative findings
(Table~\ref{tab:sample_counts}).
We do not derive this tiering from a single canonical clinical risk-stratification
instrument; rather, following the precedent established by recent conformal
selective-prediction work in clinical triage~\cite{kwon2026}, we define the tiering
directly through manual clinical-judgment grouping of SCP-ECG classes, and treat it
as a stated, documented modeling assumption rather than a claim of ground truth.
The tier boundaries are broadly consistent with how the AHA/ACC/HRS practice
standards for electrocardiographic monitoring~\cite{sandau2017} and the AHA/ACCF/HRS
recommendations for ECG standardization, Part VI~\cite{wagner2009}, stratify
analogous conditions by monitoring priority and ischemic risk, though no formal
one-to-one derivation from those instruments is performed.
Table~\ref{tab:sample_counts} documents, for every grouping, the rationale
underlying its tier assignment.

Two groupings reflect a conservative rather than a precise classification.
MI subtypes (AMI, IMI, LMI, PMI) cannot be separated into acute and chronic
presentations at this label granularity, and are treated as Critical by default
rather than assumed benign---a judgment call stated explicitly in the table.
The \texttt{\_AVB} bucket merges all degrees of AV block and cannot be disentangled
from high-grade block; it is likewise placed in Critical.

PTB-XL contains no ventricular tachycardia or ventricular fibrillation records
across any of its ten stratified folds; those rhythms are absent from the dataset
entirely.
As a consequence, atrial fibrillation (AFIB) and atrial flutter (AFLT) represent
the most severe rhythm disorders available at this granularity, and they are assigned
to the Critical tier rather than Important.
This placement reflects a dataset constraint rather than a claim that AFib carries
the same urgency as VT/VF.
``Critical'' in this framework means \emph{most severe available in PTB-XL}, not
ICU-grade emergency; we state this distinction explicitly so that any replication on
a dataset that does include VT/VF can lower AFIB and AFLT to Important without
contradicting our tier rationale.

% ─────────────────────────────────────────────────────────────────────────────
\subsubsection{Calibration procedure}
\label{sec:calibration}

Given the tier assignment, each tier is mapped to a target per-class
false-negative rate $\alpha_\mathrm{tier}$ (Table~\ref{tab:sample_counts}), and
class-conditional calibration is performed independently per class following
Algorithm~\ref{alg:aw-ccrc}.
Calibration uses the Learn-Then-Test (LTT) framework~\cite{angelopoulos2021ltt}
with an exact one-sided binomial risk test.
For each class $k$ and candidate threshold $\lambda$ on a grid $\Lambda$ of
$G{=}1000$ equally-spaced points in $[0,1]$, the null hypothesis
$H_0\colon\mathrm{FNR}(\lambda)\ge\alpha_k$ is tested using the binomial cumulative
distribution function evaluated at the observed miss count $k_g$ among the $n_k^+$
calibration positives:
\[
  p_g = \Pr\!\bigl[X \le k_g \;\big|\; X \sim \mathrm{Bin}(n_k^+,\,\alpha_k)\bigr].
\]
Working from high $\lambda$ downward and exploiting the monotonicity of
$\mathrm{FNR}(\lambda)$, the search stops at the first $\lambda$ that rejects $H_0$
at level $\delta_k$.
The resulting per-class threshold $\lambda_c^{(k)}$ defines class $k$'s contribution
to the multi-label prediction set
$\mathcal{C}(x) = \{k\in\mathcal{Y} :
  \lambda_c^{(k)}\neq\varnothing,\;f_\theta(x)_k\ge\lambda_c^{(k)}\}$.

To provide a simultaneous guarantee across all supported classes, we apply
Bonferroni correction: $\delta_k = \delta\,/\,|\mathcal{Y}_\mathrm{sup}|$ with
$\delta{=}0.10$ and $|\mathcal{Y}_\mathrm{sup}|{=}24$, giving
$\delta_k\approx 0.0042$.
This between-class correction is integral to the procedure; omitting it would reduce
the guarantee to a per-class statement at the full $\delta$ rather than the stated
simultaneous bound.

A class is certifiable only if its calibration-fold positive count $n_k^+$ satisfies
$(1-\alpha_k)^{n_k^+}\le\delta_k$---the condition under which even flagging every
positive ($k{=}0$ misses at $\lambda{=}0$) would reject $H_0$ at $\delta_k$.
This LTT certification floor substantially exceeds the weaker quantile-existence
floor $n_k^+\ge\lceil 1/\alpha_k\rceil-1$: with $\delta{=}0.10/24$, the minimum
certifiable $n_k^+$ is approximately~272 for the Critical tier ($\alpha{=}0.02$),
107 for the Important tier ($\alpha{=}0.05$), and~53 for the Benign tier
($\alpha{=}0.10$).
Table~\ref{tab:sample_counts} reports positive-example counts against the
tier-specific quantile-existence floors ($n_k^+\ge 49\,/\,19\,/\,9$ for Critical /
Important / Benign); classes that exceed this weaker floor but fall short of the LTT
certification floor are flagged with a superscript and reported as uncertifiable in
Table~\ref{tab:primary_calibration}.
AFIB ($n{=}151$) and \texttt{\_AVB} ($n{=}83$) are the primary examples: both exceed
the quantile-existence floor for Critical ($n\ge 49$) and therefore appear as
Supported in Table~\ref{tab:sample_counts}, but neither reaches the LTT
certification floor of~272 and both are therefore uncertifiable in
Table~\ref{tab:primary_calibration}.

Two classes exceed a naive fixed threshold of $n\ge 10$ yet fall below the exact
LTT requirement for their assigned tier: LMI ($n{=}20$ against a Critical-tier
requirement of~49) and SVARR ($n{=}15$ against an Important-tier requirement
of~19).
One Benign-tier class, RAO/RAE ($n{=}10$), meets its quantile-existence requirement
($n\ge 9$) by a margin of exactly one example.
These classes remain in the calibration procedure, but their resulting thresholds
are reported in subsequent tables with an explicit finite-sample caveat rather than
presented as equivalent in reliability to classes with substantially larger
calibration support.

PMI is excluded from calibration entirely.
With only $n{=}2$ positive calibration examples, the best guarantee achievable
under Algorithm~\ref{alg:aw-ccrc} is $\alpha_\mathrm{min}\approx 0.33$, looser
than every tier target in this scheme, including the most permissive
(Benign, $\alpha{=}0.10$).
No re-assignment of PMI's target $\alpha$ within the tier structure resolves this,
since the limitation is one of available sample size rather than tier placement.
We therefore omit PMI from class-conditional calibration and from per-class
false-negative-rate reporting in Section~\ref{sec:results}, and report this
omission explicitly rather than present a threshold the calibration data cannot
support.

% ─────────────────────────────────────────────────────────────────────────────
\subsubsection{Sensitivity analysis}
\label{sec:sensitivity}

% Unchanged from draft — reproduced for completeness.
Because the numeric $\alpha$ values attached to each tier are a modeling choice
rather than an externally validated quantity, we repeat the full calibration and
evaluation procedure under two additional configurations---a stricter mapping
($\alpha = 0.01/0.03/0.08$ for Critical / Important / Benign) and a looser one
($\alpha = 0.05/0.08/0.15$)---alongside the primary configuration
(0.02/0.05/0.10).
We report per-class false-negative-rate ordering, the risk--coverage ordering across
methods, and the operational flag rate for all three configurations side by side
(Section~\ref{sec:results}).
Findings that hold across all three configurations are reported as primary results;
any metric sensitive to the specific tier-to-$\alpha$ mapping is reported explicitly
as such, rather than presented only under the configuration most favorable to the
method.
